\subsection{A First-Order Differential Dynamic Programming Scheme}
\label{subsec:userddp}

The same problem can be attacked by a first-order variant of differential
dynamic programming: rather than a Riccati sweep carrying both derivatives of
the cost-to-go, it carries only the first, and solves one small program per
period sweeping forward.  Write $k$ for the sweep index.  Here
$B[t]$ is the energy \emph{after} period $t$, with $B[0] = B_0$, so $B[T]$ is
the terminal energy $B^{T+1}$ of \eqref{eq:cp_all}.

\subsubsection{The subproblem at stage $t$}
Sweep $k$ solves, for $t = 1, \dots, T$ in order,
\begin{align}
\min_{P_{\text{Subs}}[t],\, P_B[t],\, B[t],\, P_B[t\!+\!1]} \;\;
  & c^t P_{\text{Subs}}[t]\,\Delta t + C_B \bigl(P_B[t]\bigr)^2 \Delta t
\nonumber \\[-2pt]
  & +\; \mu^{k-1}[t\!+\!1] \Bigl( B^{k-1}[t\!+\!1] - B[t]
        + \Delta t\, P_B[t\!+\!1] \Bigr)
\label{eq:ddp_obj}
\end{align}
subject to the period's own balance, dynamics and limits,
\begin{alignat}{2}
& P_{\text{Subs}}[t] + P_B[t] = p_L^t
\label{eq:ddp_bal} \\
& B[t] - B^{k}[t\!-\!1] + \Delta t\, P_B[t] = 0
\label{eq:ddp_dyn} \\
& B^{\min} \leq B[t] \leq B^{\max}, \quad
  P_B^{\min} \leq P_B[t] \leq P_B^{\max} .
\label{eq:ddp_lim}
\end{alignat}
Two couplings appear, and they are not symmetric.  The state entering
\eqref{eq:ddp_dyn} is $B^{k}[t\!-\!1]$, produced by the \emph{current} sweep a
moment earlier, and it is imposed as a hard constraint.  The value of the future
enters only through the last term of \eqref{eq:ddp_obj}, and both $\mu$ and $B$
there carry the superscript $k\!-\!1$: they are numbers left over from the
\emph{previous} sweep.

\subsubsection{What is carried between sweeps}
After solving \eqref{eq:ddp_obj}--\eqref{eq:ddp_lim} the scheme reads the
multiplier of the period's own dynamics,
\begin{equation}
\mu^{k}[t] = -\,\text{dual}\bigl(\eqref{eq:ddp_dyn}\bigr),
\label{eq:ddp_mu}
\end{equation}
which is the marginal value of one more unit of stored energy at $t$.  No
separate dual update is performed: \eqref{eq:ddp_mu} is the whole of the
feedback.  Sweeps repeat until the trajectories settle,
\begin{equation}
\max\Bigl(
  \bigl\| B^{k} - B^{k-1} \bigr\|_2,\;
  \bigl\| \mu^{k} - \mu^{k-1} \bigr\|_2
\Bigr) < \epsilon .
\label{eq:ddp_conv}
\end{equation}

\subsubsection{Correspondence, and the one term that is missing}
$\mu[t]$ is the same object as the value gradient $V_B^t$ of
Section~\ref{subsec:kkt}: both are the marginal worth of stored energy at $t$,
and \eqref{eq:kkt_costate} is the statement that a converged sweep would satisfy.
The box multipliers of \eqref{eq:ddp_lim} correspond to
$\lambda_{B\min}^t, \lambda_{B\max}^t$ and to $\nu_{\min}^t, \nu_{\max}^t$
exactly as before.  The two schemes are therefore carrying the same quantities.

What has no counterpart is $V_{BB}^t$.  The coupling term in
\eqref{eq:ddp_obj} is \emph{linear} in $B[t]$: it carries the slope of the
cost-to-go evaluated at the previous sweep's state and nothing more.  The true
cost-to-go has curvature, so as $B[t]$ moves during a sweep the marginal value of
stored energy moves with it, and a frozen slope cannot follow.  Two consequences
follow, both measured below: the optimum is not quite a fixed point of the
sweep, and the iteration does not settle onto it.

\subsubsection{What it does on this instance}
\label{subsec:ddp_results}
On the $T = 6$ instance of Table~\ref{tab:cp_data} with all bounds imposed,
the scheme is implemented exactly as
\eqref{eq:ddp_obj}--\eqref{eq:ddp_conv} and started from
$B^0[t] = B_0$, $\mu^0[t] = 0$.

\emph{The subproblem is right.}  Supplied with the true multipliers $\mu^\star$
and states $B^\star$ taken from a centralized solve, a \emph{single} sweep
reproduces the centralized optimum to $1.7\times10^{-8}$ in objective and
$6\times10^{-7}$ in $P_B$.  The scheme's fixed point is the correct answer.

\emph{The iteration does not reach it.}  From the cold start above it settles
into a period-two limit cycle: two objectives alternating indefinitely, the
state residual of \eqref{eq:ddp_conv} pinned at $2.1\times10^{-1}$, still
cycling after $300$ sweeps, $2.1\times10^{-3}$ short in objective and up to
$121$~kW out in dispatch.  Damping the update,
$x \leftarrow (1-a)\,x^{k-1} + a\,x^{k}$, shrinks the cycle without closing it:

\begin{center}
\footnotesize
\begin{tabular}{ccc}
\toprule
$a$ & residual after $2000$ sweeps & gap in objective \\
\midrule
$1.0$ & $2.1\times10^{-1}$ & $2.1\times10^{-3}$ \\
$0.5$ & $2.0\times10^{-2}$ & $2.9\times10^{-4}$ \\
$0.1$ & $6.1\times10^{-4}$ & $3.2\times10^{-5}$ \\
\bottomrule
\end{tabular}
\end{center}

\noindent
The reason is visible in the oracle test: started \emph{at} the optimum, the
sweep returns multipliers differing from $\mu^\star$ by $9.2\times10^{-4}$.  The
optimum is very nearly, but not exactly, a fixed point of the linearised map, so
damping only shrinks the neighbourhood the iterate circles in.  For contrast,
the second-order method of Section~\ref{subsec:kkt} reaches
$2.2\times10^{-10}$ on the identical instance in $16$ iterations.
